% Transcription — fonds Alexandre Grothendieck, shelfmark 151, batch 3 (pages 41-60).
% Established from the facsimile archives/batches/151/batch-03.pdf (a mirror of
% https://grothendieck.umontpellier.fr/151.pdf), per the skill
% transcribe-grothendieck.
%
% Every one of the twenty pages carries the mathematics; none is skipped. The
% text « Structures stratifiées » runs on without a break from batch 2, which
% stopped mid-sentence at the foot of page 40; page 41 completes that sentence.
% The author numbers his sheets 9 to 18 at the head of the odd archivist pages
% (42, 44, 46, ..., 60), each sheet written on both sides — the continuation of
% the series 1 to 8 that batch 2 carried.
%
% The batch contains three movements. Pages 41-43 close the discussion of
% gerbes on the sphere begun in batch 2 and end on the programme this whole
% text serves: the homotopy type of a stratified space whose pieces are
% K(pi,1)'s. Pages 44-55 are section 2, « Stratifications globales :
% formalisme 1/2 simplicial (sans tubes) », which builds the semi-simplicial
% space X_{Delta_*} attached to a locally finite family of closed subspaces
% and asks when X can be recovered from it. Pages 56-60 open section 3,
% « Stratifications globales : introduction des tubes », whose equation
% numbering restarts at (1).
%
% The author writes « 1/2 simplicial » as a fraction on pages 44, 47 and 50,
% and spells « semi-simpliciale » out on pages 48 and 55; both are kept as
% written. He also writes « Quand aux » for « Quant aux » on page 49 — kept,
% unflagged, since it is orthography and not mathematics.
%
% On uncertainty: \uncertain{} marks a reading that carries mathematical weight
% and that we do not hold to be sure; \ill marks what we do not read at all.
% Function words whose sense is forced by the sentence — a « où », a « donc » —
% are given plainly. Pages 54 and 58 carry long marginal notes written
% diagonally, several times reworked; we give what reads and say so. Pages 58
% and 60 are written over so heavily in places that we transcribe the skeleton
% of the argument and mark the compression.
%
% The dating is the inventory's own and spans the group's range; the notes
% themselves carry no date.
%
% Pass: Opus 5 (claude-opus-5), 2026-08-15 — first pass, unchecked against the pages by a human.
\documentclass[11pt,a4paper]{article}
% Le préambule commun à toutes les transcriptions du fonds Grothendieck.
%
% Il tient en une idée : **l'incertitude se compose, elle ne se résout pas.**
% Une transcription qui lisse un mot douteux a détruit la seule chose pour
% laquelle on la faisait — un lecteur doit pouvoir savoir, à chaque endroit,
% ce qui était sur le feuillet et ce qui a été ajouté. Les six macros qui
% suivent sont donc l'appareil critique, et non de la décoration.
%
% Les mêmes commandes sont reconnues par `scripts/render.mjs`, qui produit la
% vue de lecture du site. Une macro ajoutée ici doit l'être aussi là-bas,
% faute de quoi le rendu échouera bruyamment — ce qui est voulu : un rendu
% silencieusement faux est pire qu'un rendu absent.

\ProvidesPackage{grothendieck}[2026/08/08 Transcriptions du fonds Grothendieck]

\usepackage{amsmath,amssymb,amsthm}
\usepackage{mathrsfs}
\usepackage{stmaryrd}
% Les diagrammes commutatifs sont partout dans ce fonds : c'est la langue
% même de la géométrie algébrique de Grothendieck, et une transcription qui
% les rendrait en prose ne transcrirait rien.
\usepackage{tikz-cd}
% Français par défaut — c'est la langue des feuillets — mais l'anglais est
% chargé aussi : la traduction et le résumé sont des documents anglais, et la
% typographie française (espaces avant les deux-points) y serait une faute.
% Ces documents basculent par \selectlanguage{english} en tête de corps.
\usepackage[english,french]{babel}
\usepackage{fontspec}
\usepackage{geometry}
\geometry{a4paper,margin=25mm,marginparwidth=28mm}
\usepackage{marginnote}
\usepackage{xcolor}
% ulem, not soul. `soul`'s \st and \ul are built on a character-by-character
% scan that XeLaTeX's font machinery breaks: the document fails with an
% undefined control sequence rather than degrading. ulem does the same two jobs
% and is Unicode-engine safe. `normalem` stops it hijacking \emph, which the
% transcriptions use for Grothendieck's own emphasis.
\usepackage[normalem]{ulem}
\usepackage{hyperref}
\hypersetup{colorlinks=true,linkcolor=black,urlcolor=black,pdfborder={0 0 0}}

% --- Métadonnées du lot -----------------------------------------------------
% Reprises telles quelles par le script de rendu, qui les affiche en tête de
% la vue de lecture. Elles fixent ce que le fichier prétend couvrir : sans
% elles, une transcription flotte sans point d'attache dans un fonds de seize
% mille feuillets.

\newcommand{\folder}[1]{\gdef\grothFolder{#1}}
\newcommand{\batch}[1]{\gdef\grothBatch{#1}}
\newcommand{\leaves}[2]{\gdef\grothFirst{#1}\gdef\grothLast{#2}}
\newcommand{\foldertitle}[1]{\gdef\grothFoldertitle{#1}}
\gdef\grothFolder{?}\gdef\grothBatch{?}\gdef\grothFirst{?}\gdef\grothLast{?}\gdef\grothFoldertitle{}

% --- Le résumé --------------------------------------------------------------
%
% Il ouvre la lecture modernisée, sous le titre « Résumé », et il est composé
% en petit. Ce n'est pas un ornement : c'est le seul passage du document qui
% ne soit pas des mathématiques, et un lecteur doit pouvoir le reconnaître
% avant de l'avoir lu — puis le sauter s'il connaît déjà le sujet, ou s'y
% arrêter s'il ne le connaît pas. Le filet qui le suit marque la reprise du
% corps.

\newenvironment{resume}
  {\small\setlength{\parskip}{0.45em}}
  {\par\medskip\hrule\medskip}

% --- Mots-clés ---------------------------------------------------------------
%
% En anglais, en clôture du résumé de la lecture modernisée : le vocabulaire
% moderne sous lequel chercher ce que les feuillets construisent. Le manifeste
% du site (`npm run manifest`) les extrait pour étiqueter la cote entière —
% cette ligne est la source unique des tags, il n'y a pas de fichier à côté.
\newcommand{\keywords}[1]{\par\smallskip\noindent
  {\footnotesize\textsc{Keywords} --- \textit{#1}}\par}

% --- L'appareil critique ----------------------------------------------------

% Le numéro de feuillet, en marge. C'est l'ancre qui permet de régler un
% désaccord sur une formule en regardant *un* feuillet plutôt qu'une chemise
% de six cents. Le site s'en sert aussi pour tourner les pages du fac-similé
% quand on fait défiler la transcription.
\newcommand{\leaf}[1]{%
  \marginnote{\footnotesize\color{gray!70}\textsf{#1}}%
  \label{leaf:#1}\ignorespaces}

% Illisible : on ne devine pas. Un mot inventé qui se lit comme les autres est
% le pire résultat possible.
\newcommand{\ill}{\textcolor{red!60!black}{[\,\ldots\,]}}

% Lecture douteuse : le mot est donné, et signalé comme incertain.
\newcommand{\uncertain}[1]{\uline{#1}}

% Ajout de l'éditeur — une lettre manquante, un mot sous-entendu.
\newcommand{\add}[1]{\textcolor{blue!50!black}{[#1]}}

% Biffé par Grothendieck lui-même. Ce qu'il a rayé fait partie du document :
% une rature dit souvent où la pensée a changé de direction.
\newcommand{\struck}[1]{\sout{#1}}

% Note du transcripteur, sur l'état matériel du feuillet.
\newcommand{\note}[1]{\par\smallskip{\small\color{gray!60!black}\textit{[#1]}}\par\smallskip}

% Note marginale de Grothendieck, distincte du corps du texte.
\newcommand{\marginal}[1]{\par\smallskip\noindent\hspace{1em}%
  {\small\color{gray!60!black}#1}\par\smallskip}

% --- Titre ------------------------------------------------------------------

\AtBeginDocument{%
  \begin{center}
    {\large\bfseries Fonds Alexandre Grothendieck --- cote n\textsuperscript{o}~\grothFolder}\\[2pt]
    {\itshape\grothFoldertitle}\\[4pt]
    {\small feuillets \grothFirst--\grothLast\ (lot \grothBatch)}\\[2pt]
    {\footnotesize\color{gray!60!black}Transcription établie d'après le fac-similé numérisé
     par l'Université de Montpellier.\\
     Les lectures incertaines sont soulignées, les illisibles notées [\,\ldots\,],
     les ajouts entre crochets.}
  \end{center}
  \vspace{1em}}

\endinput


\folder{151}
\batch{3}
\pages{41}{60}
\dating{[à partir de 1981-à partir de 1990]}
\watermark{Édition de démonstration}
\foldertitle{Topos stratifié (1981) : notes manuscrites (s.d.).}

\begin{document}

\section*{Gerbes sur une somme amalgamée — fin}

\page{41}
\ldots uns « évidents », qui ici dit que la donnée d'une gerbe $\mathcal{G}_X$
sur $X$ revient : celle des gerbes $\mathcal{G}_{X^{*}}$ sur $X^{*}$ et
$\mathcal{G}_{Y,X}$ sur $\mathcal{V}_{Y,X}$, et d'une \struck{iso.}
équivalence de \struck{telles} gerbes

\[
  (24) \qquad
  \mathcal{G}_{Y,X}\bigl|\,\mathcal{V}^{*}_{Y,X}
  \;\simeq\;
  \mathcal{G}_{X^{*}}\bigl|\,\mathcal{V}^{*}_{Y,X} .
\]

Or sur un $K(\pi,1)$, une $A$-gerbe s'interprète canoniquement comme une
extension de $\pi$ par $A$. Donc la donnée d'une gerbe sur $X$
\struck{$S^{2}$} revient : celle \struck{\ill} de deux extensions, l'une de
$\pi_1(Y)$ par $A$, l'autre de $\pi_1(X^{*})$ par $A$, et d'un isom. entre les
deux extensions de $D$ par $A$ qui s'en déduisent \ill\ (une description
modulo \ill\ que les « pièces » $Y$, $X^{*}$, $\mathcal{V}^{*}_{Y,X}$ sont
connexes, et que leur $\pi_2$ est nul).

\marginal{$D \to \pi_1(Y)$, $D \to \pi_1(X^{*})$.}

\page{42}
Dans le cas élémentaire de $S^{2}$ \add{de (23)}, on trouve donc que la donnée
d'une gerbe sur $S^{2}$ équivaut : celle d'un automorphisme de l'extension
triviale de $\mathbb{Z}$ par $A$, donc : un élément de
$\mathrm{Hom}(\mathbb{Z}, A) \simeq A$, d'où la valeur correcte

\[
  (25) \qquad H^{2}(S^{2}, A) \;\simeq\; A ,
\]

un magnifique succès ! Sans aller plus loin dans \struck{des} les
vérifications, c'est un \ill\ que l'intuition qu'il y a une théorie des
\struck{limites inductives} $\varinjlim$ de types d'homotopie (ou un
théorème ; le cas \uncertain{Kampen}, disant que le type d'homotopie d'une
$\varinjlim$ de topos (topos « usuels ») est la $\varinjlim$ des types
d'homotopie), est bien correcte. En principe, pour

\page{43}
des « espaces stratifiés » (en un sens qui va être précisé plus bas) dont les
« pièces » sont des $K(\pi,1)$, tels les espaces modulaires de Teichmüller et
leurs compactifiés, et leurs sous-espaces et voisinages tubulaires \ill,
qu'on peut leur \uncertain{associer} \uncertain{quasi} : leur stratification
canonique, le type d'homotopie doit pouvoir s'exprimer en termes d'un système
inductif convenable (le plus souvent : ens. d'indices fini) de groupoïdes
fondamentaux des dites pièces.

Il me semble d'ailleurs, réflexion faite, que cette théorie « d'intégration »
des types d'homotopie, que naguère je voyais comme un lien au grand yoga des
relations entre structures homotopiques (ou catégories, sur un topos disons)
et les \ill\ groupoïdes, est en fait relativement triviale, et justiciable
sans doute de techniques d'explicitation des plus standard. Je me réserve d'y
revenir là-dessus par la suite.

\section*{2. Stratifications globales : formalisme $\tfrac{1}{2}$ simplicial (sans tubes)}

\page{44}
Pour simplifier, je vais me placer sur un espace topologique $X$ — par la
suite $X$ deviendra un topos quelconque. Les constructions qui suivent,
relatives à une « stratification globale », seront « faisceautisées » sur $X$
par le \uncertain{lecteur}, de la façon habituelle — ce qui amènera \add{alors}
à imposer des conditions supplémentaires de connexité et de locale connexité,
qui pour le moment ne sont pas nécessaires. De même, il n'y a pas lieu pour le
moment de s'encombrer d'hypothèses de locale trivialité usuelle des
plongements d'espaces.

Soit $I$ un ensemble ordonné, \struck{\ill}

\[
  (1) \qquad (X_i)_{i \in I}, \qquad X_i \subset X
\]

une famille de \struck{parties} sous-espaces de $X$. On suppose

\[
  (1') \qquad
  \begin{cases}
    \text{a) Les } X_i \text{ fermés, la famille } (X_i) \text{ localement finie.} \\
    \text{b) } i \leq j \implies X_i \subset X_j .
  \end{cases}
\]

Posant

\[
  (2) \qquad
  X_{\Delta_0} \;=\; \coprod_{i \in I} X_i
  \;=\; \bigl\{ (x,i) \in X \times I \ \big|\ x \in X_i \bigr\}
\]

on a un morphisme canonique

\page{45}
\[
  (3) \qquad X_{\Delta_0} \longrightarrow X
\]

et l'hypothèse a) signifie que ce morphisme \struck{\ill} est \emph{fini} —
i.e. propre séparé et à fibres finies. C'est aussi une immersion locale.

On introduit une partie fermée

\[
  (4) \qquad
  X_{\Delta_1} \;\subset\; X_{\Delta_0} \times X_{\Delta_0}
  \;=\; \coprod_{(i,j) \in I \times I} X_i \times X_j ,
\]
\[
  \phantom{(4)} \qquad
  X_{\Delta_1} \underset{\text{déf}}{=} \coprod_{i \leq j} \Gamma_{i,j}
\]

\[
  (5) \qquad
  \Gamma_{ij} \subset X_i \times X_j \quad
  \text{« graphe de l'inclusion } X_i \hookrightarrow X_j \text{ ».}
\]

On voit alors que les deux projections du diagramme (6)

\begin{tikzcd}[column sep=large, row sep=small]
& X_{\Delta_1} \arrow[dl, "\sigma"'] \arrow[dr, "b"] & \\
X_{\Delta_0} & & X_{\Delta_0}
\end{tikzcd}

ont respectivement les propriétés suivantes :

\[
  (7) \qquad
  \begin{cases}
    \text{a) } X_{\Delta_1} \xrightarrow{\ \sigma\ } X_{\Delta_0}
      \text{ est un revêtement étale} \\
    \qquad (\text{sur chaque pièce } X_i)
      \text{ — en fait constant,} \\
    \qquad \text{de fibre égale à } I_i = \{ j \in I \mid j \geq i \}, \\
    \qquad \text{donc rev. fini si les } I_i \text{ sont finis} \\
    \text{b) } X_{\Delta_1} \xrightarrow{\ b\ } X_{\Delta_0}
      \text{ est une immersion locale.}
  \end{cases}
\]

\page{46}
Par ailleurs $X_{\Delta_1}$ est le graphe d'une relation d'ordre sur
$X_{\Delta_0} = \{ (x,i) \in X \times I \mid x \in X_i \}$, avec

\[
  (8) \qquad (x,i) \leq (y,j) \iff x = y, \quad i \leq j
\]

i.e. c'est la relation d'ordre induite par l'ordre produit de $X \times I$,
$X$ étant \struck{\ill} muni de l'ordre discret. Par suite, le couple
$\bigl( X_{\Delta_0},\ X_{\Delta_1} \subset X_{\Delta_0} \times X_{\Delta_0}
\bigr)$ peut être regardé comme un « objet ordonné de la catégorie des espaces
topologiques » — d'où une catégorie dans la catégorie des espaces
topologiques, ayant $X_{\Delta_0}$ comme objet des objets, $X_{\Delta_1}$
comme objet des flèches, $\sigma$ et $b$ de (6) comme morphismes source et
but, et la composition des flèches étant le morphisme (9)

\begin{tikzcd}[column sep=large, row sep=small, nodes={font=\scriptsize}]
X_{\Delta_2} \underset{\text{déf}}{=} (X_{\Delta_1}, b) \times_{X_{\Delta_0}} (X_{\Delta_1}, \sigma) \arrow[r] \arrow[d, hook] & X_{\Delta_1} \arrow[d, hook] \\
X_{\Delta_0} \times X_{\Delta_0} \times X_{\Delta_0} & X_{\Delta_0} \times X_{\Delta_0}
\end{tikzcd}

induit par la projection $(x,y,z) \mapsto (x,z)$. Plus généralement, on
introduit en

\page{47}
termes de (6) les produits fibrés multiples de $X_{\Delta_1}$ sur
$X_{\Delta_0}$, pour $r \in \mathbb{N}$

\[
  (10) \qquad
  \begin{cases}
    X_{\Delta_r} = \bigl\{ (x_0, \ldots, x_r) \in X_{\Delta_0}^{\,r+1}
      \ \big|\ x_0 \leq \cdots x_r \bigr\} \\[2pt]
    \phantom{X_{\Delta_r}} \simeq \bigl\{ (x, i_0, \ldots, i_r)
      \in X \times I^{\,r+1} \ \big| \\
    \qquad\qquad x \in X_{i_0},\ i_0 \leq \cdots \leq i_r \bigr\} \\[2pt]
    \phantom{X_{\Delta_r}} \simeq \coprod_{i_{*} \in I(\Delta_r)} X_{i_{*}} \\[2pt]
    \phantom{X_{\Delta_r}} \simeq \coprod_{i_0 \in I} X_{i_0} \times I_i(\Delta_{r-1})
  \end{cases}
\]

où

\[
  (11) \qquad
  \begin{cases}
    I(\Delta_r) = \mathrm{Hom}_{\text{ens. ord.}}(\Delta_r, I) \\[2pt]
    X_{i_{*}} = X_{i_0, i_1, \ldots, i_r} \underset{\text{déf}}{=} X_{i_0}
      \quad \text{si } i_{*} \in I(\Delta_r) .
  \end{cases}
\]

\marginal{$I(\Delta_r)$ : ens. des drapeaux de longueur $r$ de $I$.}

Donc on peut considérer, soit
$X_{\Delta_r} \subset X_{\Delta_0}^{\,r+1}$, soit
$X_{\Delta_r} \subset X \times I^{\,r+1}$.

L'homomorphisme « de composition » (9) permet alors, de façon bien connue, de
considérer les $X_{\Delta_r}$ comme étant les composantes d'un espace
$\tfrac{1}{2}$ simplicial

\[
  (12) \qquad X_{\Delta_{*}} = (X_{\Delta_r})_{r \geq 0}
\]

d'où l'application $\tfrac{1}{2}$ simpliciale

\[
  (13) \qquad \alpha^{*} : X_{\Delta_r} \longrightarrow X_{\Delta_{r'}}
\]

associée à une application croissante

\page{48}
\[
  (14) \qquad \alpha : \Delta_{r'} \longrightarrow \Delta_r
\]

est donnée par

\[
  \alpha^{*}(x, i_{*}) = \bigl( x,\ \alpha^{*}(i_{*}) \bigr),
  \qquad \alpha^{*}(i_{*}) \underset{\text{déf}}{=} i_{*} \circ \alpha
\]

i.e. elle est induite par la structure semi-simpliciale produit dans le
deuxième membre de

\[
  (15) \qquad X_{\Delta_{*}} \subset X \times I_{\Delta_{*}} .
\]

\textbf{Remarques.} a) Toutes ces constructions se généralisent \ill\ où on
remplacerait \add{l'ens. ordonné} $I$ par une catégorie $\mathcal{C}$, et
l'application $X \mapsto X_i$ par un foncteur $\mathcal{C} \to$ catégorie des
parties fermées de $X$ et de leurs inclusions.

b) Si on part d'une famille $(X_i)_{i \in I}$, \add{$I$ un ens.} sans plus, on
peut introduire ad hoc une relation d'ordre dans $I$ par \struck{\ill}

\[
  i \leq j \iff X_i \subseteq X_j .
\]

Mais dans la présentation adoptée, il n'est pas nécessaire de supposer que
\add{(l'appl. croissante)} $I \to$ ens. des parties fermées de $X$ soit
injective, et encore moins que ce soit un plongement d'ens. ordonnés, i.e. que
l'on ait

\[
  X_i \subseteq X_j \implies i \leq j
\]

qui est une réciproque de (1' b)).

\page{49}
On a

\[
  (16) \qquad
  \begin{cases}
    \text{a) si } \alpha : \Delta_{r'} \to \Delta_r
      \text{ est tel que } \alpha(0) = 0, \text{ alors} \\
    \qquad X_{\Delta_r} \xrightarrow{\ \alpha^{*}\ } X_{\Delta_{r'}}
      \text{ est un morphisme de rev. étale} \\
    \qquad \text{(en particulier } X_{\Delta_r} \to X_{\Delta_0}, \\
    \qquad\quad (x_0, \ldots, x_r) \mapsto x_0
      \text{ ou } (x, i_0, \ldots, i_r) \mapsto x, \\
    \qquad\quad \text{est un morphisme de rev. étale).}
  \end{cases}
\]

En fait, la \struck{fibre de} restriction de $X_{\Delta_r}$ au dessus du
composant $X_{i'_{*}}$ (de $X_{\Delta_{r'}}$), où
$i'_{*} = \{ i'_0 \leq \cdots \leq i'_{r'} \}$, \struck{\ill} s'identifie à un
rev. étale constant, de fibre type l'ens. des drapeaux $i_{*}$ de $I$, de type
$\Delta_r$, qui s'envoient sur $i'_{*}$ par $\alpha^{*}$ :

\[
  (17) \qquad
  \alpha^{*-1}\bigl(
    \underbrace{X_{i'_{*}}}_{\textstyle = X_{i'_0} \times \{i'_{*}\}}
  \bigr)
  \;\simeq\;
  X_{i'_{*}} \times
  \bigl\{ i_{*} \in I_{\Delta_r} \ \big|\ \alpha^{*}(i_{*}) = i'_{*} \bigr\} .
\]

Quand aux applications semi-simpliciales générales (13) (sans condition
$\alpha(0) = 0$ sur $\alpha$) on peut dire seulement qu'elles sont des immer.
locales. En fait, tous les $X_{\Delta_r}$ sont au dessus de $X$,
\uncertain{dont}

\[
  (18) \qquad X_{\Delta_r} \longrightarrow X \quad \text{imm. locale.}
\]

\page{50}
\textbf{NB} On aurait intérêt à éliminer les \struck{\ill} « redondances », en
se bornant dans la définition de $X_{\Delta_r} \subset X \times I_{\Delta_r}$
aux drapeaux « non dégénérés » \struck{de} $\subseteq I_{\Delta_r}$, i.e. aux
suites strictement croissantes. Cela oblige alors à se restreindre aux
\add{appl.} $\tfrac{1}{2}$ simpliciales relatives : \struck{des}
$\alpha : \Delta_{r'} \to \Delta_r$ qui sont strictement croissants. La
définition de $X_{\Delta_0}$ n'est pas changée, tandis que $X_{\Delta_1}$ se
trouve remplacé par $X'_{\Delta_1}$ disjoint de
$X_{\Delta_0} \times X_{\Delta_0}$.

\medskip

\textbf{Reconstitution de $X$ à partir du syst. $\tfrac{1}{2}$ simplicial
$X_{\Delta_{*}}$.} Considérons d'abord le diagramme (19)

\begin{tikzcd}[column sep=small, row sep=small]
& X_{\Delta_1} \arrow[dl, "\sigma"'] \arrow[dr, "b"] & \\
X_{\Delta_0} \arrow[dr, dashed] & & X_{\Delta_0} \arrow[dl, dashed] \\
& X &
\end{tikzcd}

je voudrais qu'il soit cocartésien, i.e. qu'on récupère $X$ comme une somme
amalgamée

\[
  (20) \qquad X \;\simeq\; X_{\Delta_0} \amalg_{X_{\Delta_1}} X_{\Delta_0} .
\]

\page{51}
Pour ceci, il faut supposer, en plus des conditions a) b) de (1'), les
conditions

\[
  (21) \qquad
  \begin{cases}
    \text{c) } X = \bigcup X_i \\[2pt]
    \text{d) } \forall (i,j) \in I \times I, \ \text{on a}\ \
      X_i \cap X_j = \bigcup_{\substack{k \in I \\ k \leq i,\ k \leq j}} X_k ,
  \end{cases}
\]

qui \uncertain{expriment} trivialement la relation (20), \struck{\ill} de
façon « formelle » :

\medskip

\textbf{Proposition.} Sous les conditions a) b) c) d) de (1')(21),
l'application canonique

\[
  (22) \qquad
  \varinjlim X_{\Delta_{*}}
  \ \bigl( X_{\Delta_0} \leftleftarrows X_{\Delta_1} \Lleftarrow X_{\Delta_2} \cdots \bigr)
  \ \xrightarrow{\ \sim\ } \ X
\]

est un homéomorphisme.

\medskip

Cette proposition est valable aussi sur les deux variantes (drapeaux
quelconques, ou \struck{dé}généré\struck{s}) pour les \struck{espaces de}
drapeaux $X_{\Delta_r}$.

\medskip

\textbf{Corollaire.} Soit
$\bigl( X_{\Delta_0},\ X_{\Delta_1} \subset X_{\Delta_0} \times X_{\Delta_0}
\bigr)$ un espace topologique ordonné, \struck{\ill}

\note{La fin de l'énoncé du corollaire est entièrement recouverte de ratures
et de reprises ; l'hypothèse est reprise au propre en tête de la page
suivante, sous le numéro (23), et c'est cette rédaction-là que nous
transcrivons.}

\page{52}
On suppose

\[
  (23) \qquad
  \begin{cases}
    \text{a) } X_{\Delta_1} \text{ est un sous-espace top. fermé} \\
    \qquad \text{de } X_{\Delta_0} \times X_{\Delta_0}
      \text{ (avec la top. induite)} \\[2pt]
    \text{b) } \forall (i,j) \in I \times I,\
      \Gamma_{ij} \underset{\text{déf}}{=}
      X_{\Delta_1} \cap (X_i \times X_j) \\
    \qquad \text{est ou bien vide, ou bien le graphe} \\
    \qquad \text{d'une immersion fermée } X_i \hookrightarrow X_j .
  \end{cases}
\]

Donc $X_{\Delta_1} \xrightarrow{\ \sigma\ } X_{\Delta_0}$ est un morphisme de
rev. \add{étale} (constant), \struck{et alors} et
$X_{\Delta_1} \xrightarrow{\ b\ } X_{\Delta_0}$ une imm. locale \ill.

Alors la relation $(i \leq j) \underset{\text{déf}}{\iff} \Gamma_{ij} \neq
\emptyset$ est une relation d'ordre sur $I$, \struck{\ill} et les $X_i$ (avec
les $X_i \to X_j$ de b)) forment un système inductif d'espaces topologiques,
indexé par $I$, aux morphismes de transition des immer. fermées ; et la donnée
de l'espace \add{top.} ordonné $X_{\Delta_0}$, \struck{ainsi} avec la
partition ensembliste de $X_{\Delta_0}$, revenant : la donnée d'un tel syst.
inductif, \struck{indexé par un ens. ordonné $I$}.

\struck{Supposons} Supposons de plus, pour nous mettre à l'aise, que $I$ est
fini, et que les $\mathrm{Inf}$ de deux éléments y existent — alors (\ill)

\[
  X = \varinjlim_{i \in I} X_i
  \ \simeq\ X_{\Delta_0} / (X_{\Delta_1} \rightrightarrows)
\]

et les $X_i \hookrightarrow X$ sont des immer. fermées de \ill

\page{53}
\textbf{Question.} Dans 21) d), faudrait-il renforcer la condition en

\[
  (24) \qquad
  \begin{cases}
    \forall\, i,j \in I, \quad i \wedge j \underset{\text{déf}}{=}
      \mathrm{Inf}(i,j) \text{ existe, et} \\[2pt]
    X_{i \wedge j} = X_i \cap X_j
  \end{cases}
  \qquad ?
\]

Il n'est pas clair qu'il le faille, car cela éliminerait, \add{ne serait-ce}
au point de vue local, des stratifications comme celles déduites du cercle
(quatre strates $A$, $B$, $\{a\}$, $\{b\}$, avec
$A \cap B = \{a\} \cup \{b\}$) comme dessiné, ou en prenant le cône dessus,
\struck{\ill} d'origine $o$, d'où une stratification locale en $o$, à cinq
strates (dont $o$) : où deux strates connexes ont une intersection disconnexe
— faut-il l'éliminer pour autant ? \struck{\ill} \struck{\ill} Il ne semble
pas. La suite le montrera bien !

\note{Le dessin porté dans la marge est un cercle coupé en deux arcs $A$ et
$B$ par les deux points $a$ et $b$, puis le cône de sommet $o$ tracé sur ce
cercle, les deux génératrices $oa$ et $ob$ marquées.}

Pour tout $x \in X$, considérons \add{l'ens. $I_x \subset I$} \struck{des}
$i \in I$ tels que $X_i \ni x$ — il y en a un nb fini \struck{\ill} (par
(1' a))), et l'ens. de ces $I$ est filtrant décroissant par (21 d)), donc il y
a un plus petit $i \in I_x$ \struck{tel que}, soit $i(x)$ :

\[
  (25) \qquad i(x) = \text{plus petit des } i \in I
  \text{ tels que } x \in X_i .
\]

Soit d'autre part, pour tout $i \in I$

\page{54}
\[
  (26) \qquad
  \begin{cases}
    \dot{X}_i \underset{\text{déf}}{=} \bigcup_{j < i} X_j
      \quad (\text{fermé de } X_i) \\[2pt]
    X_i^{*} \underset{\text{déf}}{=} X_i \smallsetminus \dot{X}_i
      \quad (\text{ouvert de } X_i) .
  \end{cases}
\]

\marginal{NB On n'exclut pas que dans (26) des $X_i$ soient $= \emptyset$, ni
\ill\ $X_i \neq \emptyset$ \ill\ i.e. remplacer $I$ par les \ill\ telles que
$X_i^{*} \neq \emptyset$ \struck{distincts} \ill\ et les
$(X_i)_{i \in I'}$ \ill}

\note{Ce NB est écrit en diagonale dans la marge de gauche, sur une dizaine de
lignes serrées et plusieurs fois reprises ; nous n'en donnons que ce qui se
lit. Il y est question de restreindre $I$ au sous-ensemble $I'$ des indices
pour lesquels $X_i^{*} \neq \emptyset$.}

\textbf{Proposition.} Les $X_i^{*}$ ($i \in I$) sont mutuellement disjoints,
et leur réunion est $X$. Si $x \in X$, $i(x)$ est aussi l'unique $i \in I$ tel
que $x \in X_i^{*}$.

\textbf{Dém.} Soit $x \in X_i^{*} \cap X_j^{*}$ ($i,j \in I$, $i \neq j$),
\struck{alors} $x \in X_i \cap X_j$, donc $\exists\, k \leq i,j$ \struck{\ill}
avec $x \in X_k$. Si on avait $k = i$ on aurait $i \leq j$ donc $i < j$,
absurde car $x \in X_j^{*} = X_j \smallsetminus \dot{X}_j$. \struck{\ill} On a
$k < i$, donc $x \in \dot{X}_i \supset X_k$, ce qui contredit
$x \in X_i^{*}$. Donc les $X_i^{*}$ sont mutuellement disjoints. D'ailleurs
$x \in X_{i(x)}^{*}$ \ill\ trivial par définition de $i(x)$, d'où la prop.

\note{Le signe entre $\dot{X}_i$ et $X_k$ est tracé d'un seul trait oblique ;
nous lisons l'inclusion que demande l'argument.}

On posera

\[
  (27) \qquad
  \begin{cases}
    \dot{X}_{\Delta_0} = \coprod_{i \in I} \dot{X}_i \\[2pt]
    X^{*}_{\Delta_0} = X_{\Delta_0} \smallsetminus \dot{X}_{\Delta_0}
      = \coprod_{i \in I} X_i^{*}
  \end{cases}
\]

puis \struck{plus} plus généralement

\[
  (28) \qquad
  \begin{cases}
    \dot{X}_{\Delta_r} = \text{image inverse de } \dot{X}_{\Delta_0}
      \text{ par } \sigma_r : X_{\Delta_r} \to X_{\Delta_0} \\
    \qquad (\text{fermé de } X_{\Delta_r}) \\[2pt]
    X^{*}_{\Delta_r} = X_{\Delta_r} \smallsetminus \dot{X}_{\Delta_r}
      = \text{image inverse de } X^{*}_{\Delta_0} \text{ par } \sigma_r \\
    \qquad (\text{ouvert de } X_{\Delta_r}) .
  \end{cases}
\]

\page{55}
Si $\alpha : \Delta_{r'} \to \Delta_r$, $\alpha(0) = 0$, alors

\[
  (29) \qquad
  \alpha^{*} : X_{\Delta_r} \longrightarrow X_{\Delta_{r'}} \quad (\text{étale})
\]

satisfait

\[
  \alpha^{*}\bigl( X^{*}_{\Delta_{r'}} \bigr) = X^{*}_{\Delta_r},
  \qquad \text{i.e.} \qquad
  \alpha^{*}\bigl( \dot{X}_{\Delta_{r'}} \bigr) = \dot{X}_{\Delta_r} .
\]

\struck{Dans les} \textbf{NB} « Moralement », $X^{*}_{\Delta_r}$ représente
une sorte de partie « non singulière » de $X_{\Delta_r}$,
$\dot{X}_{\Delta_r}$ la partie singulière ou « bord ». On fera attention que
les $\dot{X}_{\Delta_r}$ ne forment pas un espace semi-simplicial.

Pourtant, notre propos est, dans une certaine mesure, de reconstituer $X$ à
partir des « strates \ill\ » $X_i^{*}$ dans $X$, dans l'esprit du n\textsuperscript{o} 1
(où il y a deux strates $X_0 = Y$ et $X_1 = X$, \struck{\ill}
($Y \subset X$ inclusion fermée), et on \ill\ donc $Y = Y^{*}$,
$X^{*} = X \smallsetminus Y$ comme dans les notations du loc. cit., où il y
avait lieu d'introduire $\mathcal{V}_{Y,X}$.

\section*{3. Stratifications globales : introduction des tubes}

\page{56}
On garde les notations précédentes.

Pour tout couple $(i \leq j) \in I \times I$, considérons

\[
  (1) \qquad
  \overline{\mathcal{V}}_{i,j} = \mathcal{V}_{X_i, X_j}
  \qquad (\text{voisinage tubulaire de } X_i \text{ dans } X_j)
\]

(cf \ill), et posons

\[
  (2) \qquad
  \dot{\mathcal{V}}_{i,j}
  = \overline{\mathcal{V}}_{ij} \cap \bigl( R^{*}_{i,j} \bigr)
  = \overline{\mathcal{V}}_{ij} \cap R_i
\]

\note{Le membre du milieu de (2) est précédé d'un premier essai entièrement
biffé ; le dernier membre s'arrête à $R_i$ au bord de la feuille.}

où on pose, pour $i,j \in I$

\[
  (3) \qquad
  R^{*}_{i,j} = \Biggl(\ \bigcup_{\substack{k \in I \text{ tels que} \\ X_k \not\supset X_i}} X_k \Biggr) \cap X_j .
\]

\marginal{NB On a aussi $X_k \subset X_j$ :
$(3\text{bis})$\quad
$R_{ij} = \bigcup_{\substack{k < j \\ \text{tels que } X_k \cap X_i \neq X_i}} X_k$.}

\marginal{$R_i$ partie fermée de $X$.}

On notera qu'on a \struck{\ill}

\[
  (4) \qquad
  R^{*}_{i,j} \cap X_i = R_{ij} \cap X_i = \dot{X}_i .
\]

En effet, si $X_k \not\supset X_i$, i.e. $X_k \cap X_i \subsetneq X_i$, comme
$X_k \cap X_i = \bigcup_{\ell \leq k,i} X_\ell$, on voit que
\struck{\ill} $\ell \leq k,i$ implique $X_\ell \subsetneq X_i$, d'où
$\ell < i$ et $X_\ell \subset \dot{X}_i$, donc
$X_k \cap X_i \subset \dot{X}_i$, d'où $R_i \cap X_i \subset \dot{X}_i$ —
l'inclusion en sens inverse étant triviale.

\page{57}
On posera enfin

\[
  (5) \qquad
  \mathcal{V}_{ij}
  = \overline{\mathcal{V}}_{ij} \smallsetminus \dot{\mathcal{V}}_{ij}
  = \overline{\mathcal{V}}_{ij} - \overline{\mathcal{V}}_{ij} \cap R_{ij}
\]

qui est donc un espace infinitésimal autour de

\[
  (6) \qquad \mathcal{V}_{ij} \cap X_i = X_i^{*} .
\]

Je poserai enfin

\[
  (7) \qquad
  \mathcal{V}^{*}_{ij}
  = \mathcal{V}_{ij} \smallsetminus \bigcup_{k < j} X_k \cap \overline{\mathcal{V}}_{ij}
  = \overline{\mathcal{V}}_{ij} \smallsetminus \dot{X}_j
  = \mathcal{V}_{ij} \cap X_j^{*} \subset \mathcal{V}_{i,j} .
\]

Notons que pour $i = j$, on trouve

\[
  (8) \qquad
  \overline{\mathcal{V}}_{ii} = X_i , \qquad
  \dot{\mathcal{V}}_{i,i} = \dot{X}_i , \qquad
  \mathcal{V}_{i,i} = \mathcal{V}^{*}_{i,i} = X_i^{*} .
\]

Nos « pierres » de construction, dans une première étape, vont être les
$\mathcal{V}^{*}_{i,j}$, $\mathcal{V}_{ij}$ avec les inclusions incidentes (9)

\begin{tikzcd}[column sep=large, row sep=small]
& X_i^{*} \arrow[d, "\text{« âme »}"] \\
\mathcal{V}^{*}_{ij} \arrow[r] \arrow[d] & \mathcal{V}_{i,j} \\
X_j^{*} &
\end{tikzcd}

ces pierres comprenant les $X_i^{*}$ comme

\page{58}
cas particulier, en faisant $i = j$.

À vrai dire, je vois maintenant que ces fameuses pierres telles quelles sont
\uncertain{des côtés mal taillées} — elles sont trop \uncertain{générales},
\ill\ pas assez pour contenir les morceaux

\[
  \mathcal{V}^{\,d'_{*}}_{d_{*},\, d''_{*}}
  \;\simeq\;
  \widetilde{N}_{d_{*};\,(d'_0,\, d''_0)} \bigl|\ D^{*}_{d_{*}}
\]

de (Str$_{\infty}$ p.~84) — trop en un sens, que l'on peut s'en tirer avec des
« pierres » plus petites \ill\ certaines des précédentes, correspondant aux
cas $i \leq j$ où $j$ est un successeur de $i$ (i.e. $\nexists\, k$ avec
$i < k < j$).

Ces pierres sont \struck{particulières} diverses \ill\ un \ill\ sens \ill\
pour des hypothèses de lissité et de locale trivialité normales (ou
« équisingularité ») convenables de type standard. Cet énoncé \ill\ réalistes
p.~ex. pour une variété alg. complexe, \add{une stratification par} des
sous-variétés alg. complexes fermées : $\mathcal{V}^{*}_{i,j}$
\struck{est} homéomorphe : un fibré sur $X_i^{*}$ (qui est une variété
topologique), à fibre une variété topologique (et
$Z^{*}_{ij} = Z_{ij}$ vis à vis \ill) : \uncertain{cylindre ouvert}
$\mathbb{R} \times C_{ij}$, où $C_{ij}$ est une variété top. compacte.

\marginal{NB Mais ce dessin (dans le domaine réel) donne une situation
\uncertain{peu} vraisemblable — puisque \ill\ « $k > V$ » \ldots\ et \ill\
devrait \ill}

\note{Cette page est entièrement reprise entre les lignes et dans les marges ;
nous en transcrivons le fil de l'argument, en marquant par \ill\ les mots
recouverts. Le NB de la marge gauche est écrit en diagonale et accompagné d'un
croquis à main levée où se lisent $X_i$, $X_i^{*}$, $R_{ij}$ et $C_{ij}$.}

(NB $\mathcal{V}_{i,j}$ serait homéomorphe : un fibré en cônes sur $C_{ij}$
\ldots). Quand $j$ n'est pas un successeur de $i$, alors $C_{ij}$ sera encore
une variété, mais non compacte. Quand

\page{59}
aux espaces \uncertain{globaux} \ill, \struck{$\mathcal{V}_{ij}$} plus
généraux que les $\mathcal{V}^{*}_{ij}$ (et les $\mathcal{V}_{ij}$), ce sont
les suivants : on considère

\[
  i \leq j \leq k \qquad \text{d'où} \qquad X_i \subset X_j \subset X_k
\]

on prend

\[
  (10) \qquad
  \mathcal{V}^{\,j}_{i,k} \underset{\text{déf}}{=}
  \overline{\mathcal{V}}_{X_i, X_k} \smallsetminus
  \bigl( R_{j,k} \cup S_{j,k} \bigr)
\]

i.e. $R_{j,k}$ est défini dans (3)(3bis), et

\[
  (11) \qquad
  S_{j,k} \underset{\text{déf}}{=}
  \bigcup_{\substack{\ell \in I \\ \text{avec } j \leq \ell \leq k}} X_\ell
\]

\marginal{NB $S_{j,k} = \emptyset$ si $j = k$ ; $S_{j,k} = X_j$ si $k$ est un
successeur \ill.}

(On a, en comparant avec (3bis), une partition de l'ens.
$I^{(k)} \underset{\text{déf}}{=} \{ \ell \in I \mid \ell \leq k \}$ en deux
parties

\[
  I^{(k)'} = \bigl\{ \ell \in I^{(k)} \ \big|\ X_\ell \not\supset X_j \bigr\}
  \quad \text{et}
\]
\[
  I^{(k)''} = \bigl\{ \ell \in I^{(k)} \ \big|\ X_\ell \supset X_j \bigr\},
\]

de sorte que

\[
  (12) \qquad
  \begin{cases}
    R_{j,k} = \bigcup_{\ell \in I^{(k)'}} X_\ell \\[2pt]
    S_{j,k} = \bigcup_{\ell \in I^{(k)''}} X_\ell
  \end{cases}
  \qquad
  \begin{aligned}
    &\text{Donc } R_{j,k} \cup S_{j,k} = \dot{X}_k \\
    &(\text{avec } R_{j,k} \cap S_{j,k} = \dot{X}_i) .
  \end{aligned}
\]

\struck{\ill} On a donc

\[
  (13) \qquad
  \mathcal{V}^{\,i}_{i,k} = \mathcal{V}^{*}_{i,k} , \qquad
  \mathcal{V}^{\,k}_{i,k} = \mathcal{V}_{i,k} ,
\]
\[
  \phantom{(13)} \qquad
  \mathcal{V}^{\,i}_{i,i} = \mathcal{V}_{i,i} = \mathcal{V}^{*}_{i,i}
  \ \ \text{d'où} \ = X_i^{*} .
\]

$\mathcal{V}^{\,j}_{i,k}$ a encore tendance à être une \ill\ homéomorphes (ou
du moins, homotopes, ou homéomorphes en un sens \struck{faible} de germes, ce
qu'il faudrait dégager \ldots) : un fibré loc. trivial de base $X_i^{*}$ —
mais comme on voit déjà dans le cas
$\mathcal{V}^{\,k}_{i,k} = \mathcal{V}_{i,k}$ : la fibre pourrait être
singulière. Dans le cas de $\mathcal{V}^{\,k}_{i,k}$,

\page{60}
avec $k$ successeur de $i$, il s'agit de singularité isolée, mais dans le cas
des $\mathcal{V}^{\,k}_{i,k}$ généraux, on peut trouver des singularités de
type arbitraire — bien sûr en fibre (prendre $X_i = X_i^{*}$ réduit à un pt !)
— ce qui diminue l'intérêt, dans le cas de stratifications
\struck{\ill} \add{ne soit pas} générales, des $\mathcal{V}^{\,j}_{i,k}$
généraux comme « pierres ». Les
$\mathcal{V}^{\,j}_{i,k} = \mathcal{V}^{*}_{i,k}$ sont déjà plus
intéressants, car \struck{\ill} c'est essentiellement un \ill\ de $X_k^{*}$,
dans une singularité — ce qui s'exprime par le fait qu'on l'obtient comme
fibré sur $X_i^{*}$ non singulier, à fibre non singulière
$\mathcal{V}^{*}_{z,\, Z_{i,k}}$ \ldots

Donc nous allons dans la suite \uncertain{systématiquement} essayer
d'exprimer les \ill\ avec les $X_i^{*}$ et les \struck{$\mathcal{V}_{ij}$}
$\mathcal{V}^{*}_{i,j}$, $\mathcal{V}_{i,j}$ « élémentaires » ($j$ successeur
de $i$), puis voir si certaines expressions peuvent prendre une forme plus
sympathique en \struck{les $X_i^{*}$ et} utilisant aussi les
$\mathcal{V}_{i,j}$, $\mathcal{V}^{*}_{i,j}$ \ill\ $i \leq j$ \ill\ plus — et
dans un \struck{dernier} temps seulement, peut-être (eu égard aux cas des
diviseurs : croisements normaux) voir comment les $\mathcal{V}^{\,j}_{i,k}$
s'insèrent dans ce formalisme.

\end{document}
